\documentclass[11pt,a4paper]{report}

\usepackage{microtype}
\usepackage{xcolor}
\usepackage{subcaption}
\usepackage{amsmath}

% Editing macros — delete for the final draft
% (\fc, \ct, \app, \sr are already defined below in this file)
\newcommand{\fig}[1]{\textcolor{purple}{[FIG: #1]}}
\newcommand{\ai}[1]{\textcolor{brown}{[AI: #1]}}

\begin{document}
% 1. FACT CHECK (Critical): Red commands immediate attention for potential inaccuracies.
\newcommand{\fc}[1]{\textcolor{red}{[FactC: #1]}} 
% 2. CITE (High): Orange highlights missing evidence or attribution.
\newcommand{\ct}[1]{\textcolor{orange}{[CITE: #1]}} 
% 3. APPENDIX (Medium): Blue indicates structural moves (shifting text out of the main body).
\newcommand{\app}[1]{\textcolor{blue}{[APP: #1]}} 
% 4. SMALL REWORD (Low): Teal is for minor stylistic edits, mirroring your friend's setup.
\newcommand{\sr}[1]{\textcolor{teal}{[Sreword: #1]}}
%=======================================================================================
\chapter{Introduction}
\section{Motivation}
% Forestry problem and wider importance
A Forest Digital Twin is a wider long-term goal for the forestry sector: a dynamic digital representation of forest systems that can support monitoring, forecasting and management planning. It is a shift from static, periodic assessment toward adaptive, data-driven planning. Rather than relying on inventories and lookup tables, Digital Twin approaches combine repeated observations, simulation, artificial intelligence(AI) and machine-learning(ML) methods to represent forests as changing systems \cite{baskent2026Review_C1_SS}. This is increasingly vital as managers face climate-driven risks—such as extreme storms, destructive winds, and severe summer heatwaves that trigger wildfires and alter growth patterns. At the same time, global forestry is shifting toward continuous cover management, demanding more responsive, stand-specific tools.\cite{ForestResearch2009Evidence_C1_SS}

Sitka spruce (\textit{Picea sitchensis}) is Britain's primary commercial species, covering approximately 50\% of managed forests \cite{ForestResearch2026Sitka_C1_SS} and contributing roughly \pounds1.1 billion annually to Scotland's economy \cite{ScottishForestry2026Future_C1_SS, mason2011Sitka_C1_SS}. %TODO: mention aberfolye here, well studied site, mostly sitka spruce lidar scans covering 21 years
Sitka spruce (\textit{Picea sitchensis}) is Britain's primary commercial species, covering approximately 50\% of managed forests \cite{ForestResearch2026Sitka_C1_SS} and contributing roughly \pounds1.1 billion annually to Scotland's economy \cite{ScottishForestry2026Future_C1_SS, mason2011Sitka_C1_SS}. The study focuses on Aberfoyle, Scotland, a Sitka spruce-dominated region with 21 years of repeated LiDAR coverage. Featuring varied topography—from steep glens and exposed, windy ridges to sheltered valley floors—it provides an ideal landscape to evaluate localised environmental influences on growth.

% Height, LiDAR and shared growth relationship
Top height is a key measure of stand development, commonly used in forestry models to estimate downstream metrics like timber volume and site productivity \cite{ForestResearch2026Sitka_C1_SS}. Traditionally measured in the field, top height can now be accurately captured at scale using the 95th percentile of airborne LiDAR return elevations \cite{suarezminguez2008Practical_C1_SS}. For most tree species, including Sitka spruce, height development over time is modelled by the Chapman--Richards equation \cite{pommerening2015Methods_C2_CR}. This function describes a sigmoid growth curve, with its parameters fixed to represent a stand's maximum height potential and its speed of growth. A fixed parameter curve only captures the average population trend, but local trajectories can deviate from this baseline, due to factors like terrain, climate, soil, wind exposure, and active management all critically constrain or enhance a stand's actual development \cite{socha2023Higher_C1_SS, cameron2015Building_C1_SS}. This creates two related questions: whether incorporating these local environmental conditions can improve spatial height prediction, and whether they are associated with plots departing from the shared growth relationship.

This leaves two key questions: whether incorporating these local drivers improves spatial height prediction, and whether they explain why specific plots depart from standard baseline growth.

Flexible ML architectures can capture complex, non-linear interactions across age, stand structure, and local environment. While models like deep neural networks (DNNs) and gradient-boosted decision trees (XGBoost) learn these relationships directly from observation, physics-informed neural networks (PINNs) provide a hybrid approach by embedding domain knowledge directly into the training objective. DNNs are known to struggle with similar growth processes, as its susceptible to overfitting or data sparsity, creating unplausible predictions for example decreasing height over time. The PINN instead uses the CR equations and derivate terms in the loss function to act as regularisation. These soft loss constraints guide predictions toward smooth, plausible growth shapes, while incorporating environmental features allows the model to capture how local conditions alter the growth trajectory. \cite{karpatne2017Theory_C4_NN, pichler2025Inferring_C4_NN}

% Research Gap & Thesis Scope
While repeated airborne laser scanning has been applied to model forest height growth and productivity \cite{Janiec2023ALSHeightGrowth, TymianskaCzabanska2021WeatherHeightGrowth, Tompalski2015ALSProductivity}, limited evidence exists on whether Chapman Richards hybrid model improves predictions of Sitka spruce top height, and using these same models as a tool to more understand the correlations between abiotic factors and forest growth. 

Many current remote-sensing models do not account for spatial dependence between nearby-identical trees, leading to data leakage and inflated performance estimates \cite{Karasiak2022Spatial_C5_EVAL, ploton2020Spatial_C5_EVAL, lynch2025Digital_C4_NN}. Relying on black-box neural networks without explainable AI limits their use for stakeholders who require interpretable insights to support decision-making.

Therefore the focus of this dissertation targets a small component of the Digital Twin pipeline: testing whether incorporating multi-source environmental data and ecological theory improves spatial height predictions, and using attribution techniques to explain local growth deviations.


\section{Aim and Objectives}
The overall aim is to evaluate the prediction and environmental attribution of repeated LiDAR-derived Sitka spruce height trajectories.

\vspace{-0.5em}
\begin{table}[!htbp]
\caption{Research Aim and Objectives}
\label{tab:aims_objectives}
\small % Slightly larger than scriptsize, very readable
\setlength{\tabcolsep}{4pt}
\renewcommand{\arraystretch}{1.2} % Tightly packs rows but keeps text legible
\begin{tabular}{p{0.04\textwidth} p{0.9\textwidth}}
\hline
\textbf{No.} & \textbf{Objective Description} \\
\hline
\textbf{1} & Create and clean the dataset provided by Forest Research, integrating it with wind, terrain, climate, and soil data across the four- and six-survey cohorts (\ref{chap:data}). \\

\textbf{2} & Establish a growth baseline by fitting a shared Chapman-Richards (CR) growth curve, and compare it to machine learning models, DNNs, and PINNs. \\

\textbf{3} & Assess the predictive performance of top height through the following sub-objectives: \newline
  3.1) Compare the DNN and PINN against baselines (Age by Age, Linear Regression, Random Forest, XGBoost) given identical feature sets. \newline
  3.2) Evaluate generalisation under spatially disjoint and temporal hold-out validation sets. \newline
  3.3) Quantify performance changes across different environmental feature combinations. \newline
  3.4) Isolate the contribution of physics guidance by comparing different loss weights for PINN architectures. \\

\textbf{4} & Determine departures from the shared growth curve through the following sub-objectives: \newline
  4.1) Determine if environmental conditions in neural models selectively reduce prediction errors on plots with the largest baseline departures. \newline
  4.2) Identify management and abiotic environmental drivers explaining the residuals using NLME, Elastic Net, and XGBoost (\ref{chap:methods}). \\

\textbf{5} & Examine local growth trajectories and how environmental impacts vary by location through the following sub-objectives: \newline
  5.1) Fit plot-specific CR curves to define individual growth trajectories. \newline
  5.2) Identify drivers of the local growth trajectories and investigate reasons for their deviations from the shared growth curve. \\
\hline
\end{tabular}
\end{table}
\vspace{-0.5em}


%Main finding paragraph (COMPLETE AT THE END)
The main finding is: 
\TODO: leave until the end. 

%---------------------

\end{document}